\documentclass[12pt,reqno]{amsart}
%%%%%%%%%%%%%%%%%%%%%%%%%%%%%%%%%%%%%%%%%%%%%%%%%%%
%								Packages									         %
%%%%%%%%%%%%%%%%%%%%%%%%%%%%%%%%%%%%%%%%%%%%%%%%%%%
\usepackage[T1]{fontenc}
\usepackage{amsmath, amssymb, amsthm, amscd, amsfonts}
\usepackage{mathtools}
\usepackage{enumerate,multicol}
\usepackage[shortlabels]{enumitem}

\usepackage[dvipsnames]{xcolor}
\definecolor{DartmouthGreen}{HTML}{00693E}
\definecolor{DartmouthDarkGreen}{HTML}{004B23}
\definecolor{DartmouthLightGreen}{HTML}{4BAF7A}

\usepackage[
  colorlinks=true,
  hyperindex,
  linkcolor=DartmouthDarkGreen,
  citecolor=DartmouthGreen,
  urlcolor=DartmouthGreen,
  pagebackref=true
]{hyperref}

\usepackage{parskip}
\usepackage[letterpaper, margin=.5in, top=1in, bottom=1.4in, headheight=42pt, headsep=14pt, footskip=50pt]{geometry}
\linespread{1.1}

\usepackage{algorithm,algpseudocode,verbatim}
\usepackage{float,caption,comment}

\usepackage{tikz,tikz-cd}
\usetikzlibrary{graphs, graphs.standard,positioning}



\usepackage{calrsfs}
\DeclareMathAlphabet{\pazocal}{OMS}{zplm}{m}{n}
\usepackage{newcent,fouriernc}

%%%%%%%%%%%%%%%%%%%%%%%%%%%%%%%%%%%%%%%%%%%%%%%%%%%
%								Copyright 								         %
%%%%%%%%%%%%%%%%%%%%%%%%%%%%%%%%%%%%%%%%%%%%%%%%%%%
\usepackage{ccicons}
\usepackage{fancyhdr}

\pagestyle{fancy}
\fancyhf{}
\fancyfoot[R]{\footnotesize \thepage}
\renewcommand{\headrulewidth}{0pt}
\renewcommand{\footrulewidth}{0pt}

\newcommand{\originalurl}{}
\newcommand{\originalurlI}{}

\fancypagestyle{firstpage}{%
  \fancyhf{}
\fancyhead[L]{\footnotesize Dartmouth College \\ Math 121: Commutative Algebra}
\fancyhead[R]{\footnotesize \href{https://www.juliettebruce.xyz/teaching/121S26}{Spring 2026}}
\fancyfoot[L]{\footnotesize \ccbyncsa\ \ \ccCopy \ Juliette Bruce 2026.\\
Licensed under \href{https://creativecommons.org/licenses/by-nc-sa/4.0/}{Creative Commons BY-NC-SA 4.0}.\\
Adapted from \ccCopy \href{https://sites.lsa.umich.edu/kesmith/}{Karen Smith} (2019); see \href{\originalurl}{original}.}
\fancyfoot[R]{\footnotesize \thepage}
\renewcommand{\headrulewidth}{0.4pt}
\renewcommand{\footrulewidth}{0.4pt}
}

\fancypagestyle{firstpage2}{%
  \fancyhf{}
\fancyhead[L]{\footnotesize Dartmouth College \\ Math 121: Commutative Algebra}
\fancyhead[R]{\footnotesize \href{https://www.juliettebruce.xyz/teaching/121S26}{Spring 2026}}
\fancyfoot[L]{\footnotesize \ccbyncsa\ \ \ccCopy \ Juliette Bruce 2026.\\
Licensed under \href{https://creativecommons.org/licenses/by-nc-sa/4.0/}{Creative Commons BY-NC-SA 4.0}.\\
Adapted from \ccCopy \href{https://sites.lsa.umich.edu/kesmith/}{Karen Smith} (2019); see \href{\originalurl}{original} and \href{\originalurlI}{original}.}
\fancyfoot[R]{\footnotesize \thepage}
\renewcommand{\headrulewidth}{0.4pt}
\renewcommand{\footrulewidth}{0.4pt}
}


\fancypagestyle{firstpageIP}{%
  \fancyhf{}
\fancyhead[L]{\footnotesize Dartmouth College \\ Math 121: Commutative Algebra}
\fancyhead[R]{\footnotesize \href{https://www.juliettebruce.xyz/teaching/121S26}{Spring 2026}}
\fancyfoot[L]{\footnotesize \ccbyncsa\ \ \ccCopy \ Juliette Bruce 2026.\\
Licensed under \href{https://creativecommons.org/licenses/by-nc-sa/4.0/}{Creative Commons BY-NC-SA 4.0}.\\
Inspired by \ccCopy\ \href{https://sites.lsa.umich.edu/kesmith/}{Karen Smith} (2019); \href{https://sites.lsa.umich.edu/kesmith/teaching/math-614-commutative-algebra/}{see}.}
\fancyfoot[R]{\footnotesize \thepage}
\renewcommand{\headrulewidth}{0.4pt}
\renewcommand{\footrulewidth}{0.4pt}
}
%%%%%%%%%%%%%%%%%%%%%%%%%%%%%%%%%%%%%%%%%%%%%%%%%%%
%								Theorems 								         %
%%%%%%%%%%%%%%%%%%%%%%%%%%%%%%%%%%%%%%%%%%%%%%%%%%%

\newtheorem{maintheorem}{Theorem}
\newtheorem{theorem}{Theorem}
\newtheorem{lemma}[theorem]{Lemma}

\newtheorem{corollary}[theorem]{Corollary}
\newtheorem{proposition}[theorem]{Proposition}
\newtheorem{conjecture}[theorem]{Conjecture}


\theoremstyle{definition}
\newtheorem{maindefinition}[maintheorem]{Definition}
\newtheorem{definition}[theorem]{Definition}
\newtheorem{setup}[theorem]{Set-Up}
\newtheorem{notation}[theorem]{Notation}
\newtheorem{convention}[theorem]{Convention}
\newtheorem{construction}[theorem]{Construction}





\setcounter{MaxMatrixCols}{18}

%%%%%%%%%%%%%%%%%%%%%%%%%%%%%%%%%%%%%%%%%%%%%%%%%%%
%								Letters									         %
%%%%%%%%%%%%%%%%%%%%%%%%%%%%%%%%%%%%%%%%%%%%%%%%%%%


\newcommand{\CC}{\mathbb{C}}
\newcommand{\FF}{\mathbb{F}}
\newcommand{\EE}{\mathbb{E}}

\newcommand{\GG}{\mathbb{G}}
\newcommand{\HH}{\mathbb{H}}
\newcommand{\II}{\mathbb{I}}
\newcommand{\KK}{\mathbb{K}}

\newcommand{\MM}{\mathbb{M}}
\newcommand{\NN}{\mathbb{N}}
\newcommand{\PP}{\mathbb{P}}
\newcommand{\QQ}{\mathbb{Q}}
\newcommand{\RR}{\mathbb{R}}
\renewcommand{\SS}{\mathbb{S}}
\newcommand{\TT}{\mathbb{T}}
\newcommand{\VV}{\mathbb{V}}
\newcommand{\WW}{\mathbb{W}}

\newcommand{\ZZ}{\mathbb{Z}}


\newcommand{\cA}{\mathcal{A}}
\newcommand{\cB}{\mathcal{B}}
\newcommand{\cC}{\mathcal{C}}
\newcommand{\cE}{\mathcal{E}}

\newcommand{\cF}{\mathcal{F}}
\newcommand{\cG}{\mathcal{G}}

\newcommand{\cH}{\mathcal{H}}
\newcommand{\cI}{\mathcal{I}}
\newcommand{\cL}{\mathcal{L}}
\newcommand{\cM}{\mathcal{M}}
\newcommand{\cO}{\mathcal{O}}
\newcommand{\cP}{\mathcal{P}}
\newcommand{\cQ}{\mathcal{Q}}
\newcommand{\cS}{\mathcal{S}}
\newcommand{\cT}{\mathcal{T}}
\newcommand{\cR}{\mathcal{R}}
\newcommand{\cU}{\mathcal{U}}
\newcommand{\cV}{\mathcal{V}}

\newcommand{\pB}{\pazocal{B}}
\newcommand{\pC}{\pazocal{C}}
\newcommand{\pO}{\pazocal{O}}
\newcommand{\pG}{\pazocal{G}}

\newcommand{\sM}{\mathsf{M}}
\newcommand{\sN}{\mathsf{N}}

\newcommand{\sQ}{\mathsf{Q}}
\newcommand{\sP}{\mathsf{P}}

\newcommand{\fm}{\mathfrak{m}}
\newcommand{\fp}{\mathfrak{p}}
\newcommand{\fq}{\mathfrak{q}}

%%%%%%%%%%%%%%%%%%%%%%%%%%%%%%%%%%%%%%%%%%%%%%%%%%%
%								Operators									         %
%%%%%%%%%%%%%%%%%%%%%%%%%%%%%%%%%%%%%%%%%%%%%%%%%%%
\DeclareMathOperator{\Id}{Id}
\DeclareMathOperator{\ev}{ev}
\DeclareMathOperator{\ord}{ord}
%
\DeclareMathOperator{\Hom}{Hom}
\DeclareMathOperator{\End}{End}
\DeclareMathOperator{\coker}{coker}
\DeclareMathOperator{\Ann}{Ann}
\DeclareMathOperator{\img}{img}
%
\DeclareMathOperator{\Frac}{Frac}
\DeclareMathOperator{\Nil}{Nil}
\DeclareMathOperator{\red}{red}
%
\DeclareMathOperator{\Spec}{Spec}
\DeclareMathOperator{\mSpec}{mSpec}
%
\DeclareMathOperator{\SL}{SL}
%
\DeclareMathOperator{\init}{in}
\DeclareMathOperator{\lex}{lex}
\DeclareMathOperator{\grlex}{grlex}
\DeclareMathOperator{\grevlex}{grevlex}
\DeclareMathOperator{\revlex}{revlex}
\DeclareMathOperator{\LT}{LT}
\DeclareMathOperator{\LM}{LM}
\DeclareMathOperator{\LC}{LC}


%%%%%%%%%%%%%%%%%%%%%%%%%%%%%%%%%%%%%%%%%%%%%%%%%%%
%								Categories									         %
%%%%%%%%%%%%%%%%%%%%%%%%%%%%%%%%%%%%%%%%%%%%%%%%%%%

\DeclareMathOperator{\comRing}{\textbf{comRing}}
\DeclareMathOperator{\Top}{\textbf{Top}}
\DeclareMathOperator{\Set}{\textbf{Set}}
\DeclareMathOperator{\alg}{\textbf{alg}}
\DeclareMathOperator{\Mod}{\textbf{Mod}}


%%%%%%%%%%%%%%%%%%%%%%%%%%%%%%%%%%%%%%%%%%%%%%%%%%%
%								MISC									         %
%%%%%%%%%%%%%%%%%%%%%%%%%%%%%%%%%%%%%%%%%%%%%%%%%%%
\makeatletter
\newcommand*{\tp}{%
  {\mathpalette\@transpose{}}%
}
\newcommand*{\@transpose}[2]{%
  % #1: math style
  % #2: unused
  \raisebox{\depth}{$\m@th#1\intercal$}%
}
\makeatother

\newcommand{\juliette}[1]{{\color{purple} \sf  Juliette: [#1]}}
%%%%%%%%%%%%%%%%%%%%%%%%%%%%%%%%%%%%%%%%%%%%%%%%%%%
%%%%%%%%%%%%%%%%%%%%%%%%%%%%%%%%%%%%%%%%%%%%%%%%%%%
\renewcommand{\originalurl}{https://sites.lsa.umich.edu/kesmith/wp-content/uploads/sites/1309/2024/07/IntegralWorksheet.pdf}
\title{Worksheet 7.1: Integral Extension Pt. I}

%%%%%%%%%%%%%%%%%%%%%%%%%%%%%%%%%%%%%%%%%%%%%%%%%%%
%%%%%%%%%%%%%%%%%%%%%%%%%%%%%%%%%%%%%%%%%%%%%%%%%%%

%%%%%%%%%%%%%%%%%%%%%%%%%%%%%%%%%%%%%%%%%%%%%%%%%%%
%%%%%%%%%%%%%%%%%%%%%%%%%%%%%%%%%%%%%%%%%%%%%%%%%%%

\begin{document}

%%%%%%%%%%%%%%%%%%%%%%%%%%%%%%%%%%%%%%%%%%%%%%%%%%%
%%%%%%%%%%%%%%%%%%%%%%%%%%%%%%%%%%%%%%%%%%%%%%%%%%%

\maketitle
\thispagestyle{firstpage}
Throughout this course, ``ring'' means \textit{commutative} ring with unity and all ring homomorphisms preserve unity. In the previous worksheets we have repeatedly moved between rings, modules, quotients, localization, and spectra. The goal of this worksheet is to introduce a language that lets us treat many ring maps as geometric families: the language of \emph{$R$-algebras}. Once this language is in place, we will distinguish two finiteness notions that sound similar but behave very differently: finite generation as an algebra and finite generation as a module.

The main bridge between these two notions is ntegrality. Integral elements are elements of an extension ring that satisfy monic polynomial equations over the base. The monic condition is what lets one convert algebraic equations into module-theoretic finiteness. This is one of the basic mechanisms behind many finiteness theorems in commutative algebra and algebraic geometry.


Let $R$ be a ring. An \emph{$R$-algebra} is a ring $S$ together with a specified ring homomorphism $\alpha:R\longrightarrow S$. Note we do not require that $\alpha$ be injective, however, the case when $\alpha$ is an inclusion, i.e., $\KK$ into $\KK[x]$, is generally the motivating example. We often think of an $R$-algebra structure on a ring $S$ as giving us a way to multiply elements of $S$ by elements of $R$, i.e., $r \cdot s \coloneqq \alpha(r)s$. Thus every $R$-algebra $S$ is naturally an $R$-module. The key difference between an $R$-module $M$ and an $R$-algebra $S$ is that for module we may only multiple elements of $M$ by elements of $R$, but for an $R$-algebra we every element of $S$ may be multiplied together. The key analogy might be: $R$-modules are to vector spaces, as $R$-algebras are to rings. 

An $R$-algebra homomorphism $\phi:S\to T$ is a ring homomorphism compatible with the structure maps from $R$, meaning that the diagram
\[
\begin{tikzcd}[row sep=4em, column sep=4em]
S \arrow[rr, "\phi"] \arrow[dr, "\alpha_{S}"]& & T \arrow[dl, "\alpha_{T}"] \\
& T &
\end{tikzcd}
\]
commutes. Equivalently, $\phi(\alpha_S(r))=\alpha_{T}(r)$ for every $r\in R$. An \emph{$R$-subalgebra} of $S$ is a subring $S' \subset S$, which contains $\alpha(R)$. Given any set $\cS \subset S$ the \emph{$R$-algebra generated by $\cS$} is the smallest $R$-subalgebra of $S$ containing $\cS$. We often write $R[\cS]$ to denote $R$-subalgebra generated by $\cS$, because it can explicitly be written down as the set of polynomials in the elements in $\cS$ with coefficients in $R$.  

\hrulefill

\begin{enumerate}[leftmargin=*, series=problems]

\item \textbf{First Encounters with $R$-algebras.} Let $\alpha:R\to S$ be an $R$-algebra.
\begin{enumerate}
    \item Verify that the rule $r\cdot s=\alpha(r)s$ makes $S$ into an $R$-module.
    \item Give an example where $\alpha$ is not injective. In your example, identify a nonzero element of $R$ that acts as zero on $S$.
    \item Show that every quotient $R/I$ is naturally an $R$-algebra.
    \item Show that if $S$ and $T$ are $R$-algebras and $\phi:S\to T$ is an $R$-algebra homomorphism, then $\phi$ is $R$-linear for the underlying $R$-module structures.
    \item Let $\KK \subseteq \EE$ be a field extension. Explain why $\EE$ is a $\KK$-algebra. What is the underlying $\KK$-module structure?
\end{enumerate}

\item \textbf{Polynomial Rings and Generation of $R$-algebras.} Let $S$ be an $R$-algebra and let $s_1,\ldots,s_n\in S$.
\begin{enumerate}
    \item Construct an $R$-algebra homomorphism
    \[
    \ev_{s_1,\ldots,s_n}:R[x_1,\ldots,x_n]\longrightarrow S,
    \qquad
    f(x_1,\ldots,x_n)\longmapsto f(s_1,\ldots,s_n).
    \]
    \item Show that the image of this map is the smallest $R$-subalgebra of $S$ containing $s_1,\ldots,s_n$. We denote this image by $R[s_1,\ldots,s_n]$.
    \item Prove that $S=R[s_1,\ldots,s_n]$ if and only if $\ev_{s_1,\ldots,s_n}$ is surjective.
    \item Deduce that $S$ is generated as an $R$-algebra by $n$ elements if and only if $S\cong R[x_1,\ldots,x_n]/I$ as an $R$-algebra for some ideal $I \subset R[x_{1},\ldots,x_{n}]$.
\end{enumerate}

\item \textbf{Examples of $R$-algebras.} In each case, identify a finite set of algebra generators over the indicated base ring. Try to choose a minimal-looking set, though you do not need to prove minimality.
\begin{multicols}{2}
\begin{enumerate}
    \item $R[x]/\langle x^2\rangle$ over $R$.
    \item $R[x,y]/\langle xy\rangle$ over $R$.
    \item $R/I$ over $R$.
    \item $\KK[t,t^{-1}]$ over $\KK$.
    \item $\ZZ[i]$ over $\ZZ$.
    \item $\KK[x,y]/\langle y^2-x^3\rangle$ over $\KK[x]$.
    \item $\KK(t)[[t]]$ over the $\KK(t)$.
    \item $\ZZ[\sqrt{n} \; n \in \ZZ_{\geq1}]$ over $\ZZ$. 
\end{enumerate}
\end{multicols}

\end{enumerate}

\hrulefill

As discussed above an $R$-algebra $S$ is also always an $R$-module with the same $R$-multiplication. This leads to two different finiteness conditions depending on whether we wish to view $S$ as an $R$-module or as an $R$-algebra. Let $S$ be an $R$-algebra.

\begin{definition}
   We say that $S$ is \emph{algebra-finite over $R$} if $S$ is finitely generated s as an $R$-algebra.
\end{definition}

\begin{definition}
	We say that $S$ is \emph{module-finite over $R$} if $S$ is finitely generated as an $R$-module.
\end{definition}

The condition $S$ be algebra-finite over $R$ is sometimes called being \emph{finite type} over $R$, i.e. we would say $S$ is finite type over $R$ if and only if $S$ is algebra-finite over $R$. Similarly, many authors will use the phrase $S$ is \emph{finite} over $R$ or $S$ is a \emph{finite $R$-algebra} to mean $S$ if a $R$-algebra that is module-finite,. In this worksheet we will attempt to restrict to using \emph{algebra-finite} and \emph{module-finite} to hopefully minimize confusion.

The key point is that module-finite is much stronger than algebra-finite. A module-finite algebra is controlled by finitely many $R$-linear combinations. An algebra-finite algebra is controlled by finitely many elements and all polynomial expressions in them.

\hrulefill

\begin{enumerate}[leftmargin=*, resume=problems]

\item \textbf{Module-Finite Implies Algebra-Finite.} Let $S$ be an $R$-algebra and suppose that $S$ is generated as an $R$-module by $m_1,\ldots,m_n\in S$.
\begin{enumerate}
    \item Show that the $R$-subalgebra $R[m_1,\ldots,m_n]$ contains every $R$-linear combination of $m_1,\ldots,m_n$.
    \item Deduce that $S=R[m_1,\ldots,m_n]$.
    \item Conclude that every module-finite $R$-algebra is algebra-finite over $R$.
\end{enumerate}

\item \textbf{Algebra-Finite Need Not Imply Module-Finite.} Assume $R$ is nonzero.
\begin{enumerate}
    \item Show that $R[x]$ is algebra-finite over $R$.
    \item Suppose $p_1,\ldots,p_n\in R[x]$. Let $d$ be the maximum of their degrees. Show that every $R$-linear combination of the $p_i$ has degree at most $d$, unless it is zero.
    \item Deduce that $R[x]$ is not module-finite over $R$. Hint: Is $x^{d+1}$ in the $R$-span of $p_1,\ldots,p_n$?
    \item Specialize to $R=\KK$ a field and restate the result in the language of vector spaces.
\end{enumerate}

\item \textbf{Permanence of Algebra-Finiteness.} The goal of this exercise is to explore how being algebra-finite plays with many natural algebraic operations: quotienting, tensoring, localizing, and composing the underlying morphisms giving the algebra structure. Consider maps of rings:
\[
\begin{tikzcd}[row sep = .3em, column sep = 4em]
R \arrow[r, "\alpha"] & S \arrow[r, "\beta"] & T
\end{tikzcd}.
\]
The map $\alpha$ gives $S$ an $R$-algebra structure. On the other hand $T$ has two different algebra structures: we may think of $T$ as an $S$-algebra via the map $\beta$ or think of $T$ as an $R$-algebra via $\beta\circ\alpha$. 
\begin{enumerate}
    \item If $S$ is algebra-finite over $R$ and $J\subseteq S$ is an ideal, prove that $S/J$ is algebra-finite over $R$.
    \item If $S$ is algebra-finite over $R$ and $T$ is algebra-finite over $S$, prove that $T$ is algebra-finite over $R$.
    \item Give an example where $\alpha$ and $\beta$ are injective and $T$ is algebra-finite over $R$, but $S$ is not algebra-finite over $R$. Hint: Let $R=\KK$, and $T = \KK[x,y]$, what happens if $S=\KK[x,xy,xy^{2},xy^{3},\ldots]$?
    \item If $S$ is algebra-finite over $R$, and $R\to T$ is a ring map, prove that $T\otimes_R S$ is algebra-finite over $T$. What are algebra generators?
    \item If $U\subset R$ is multiplicatively closed and $S$ is algebra-finite over $R$, prove that $U^{-1}S$ is algebra-finite over $U^{-1}R$.
\end{enumerate}

\item \textbf{Permanence of Module-Finiteness.} The goal of this exercise is to explore how being module-finite plays with many natural algebraic operations: quotienting, tensoring, localizing, and composing the underlying morphisms giving the algebra structure. Consider maps of rings:
\[
\begin{tikzcd}[row sep = .3em, column sep = 4em]
R \arrow[r, "\alpha"] & S \arrow[r, "\beta"] & T
\end{tikzcd}.
\]
The map $\alpha$ gives $S$ an $R$-module structure. On the other hand $T$ has two different module structures: we may think of $T$ as an $S$-module via the map $\beta$ or think of $T$ as an $R$-module via $\beta\circ\alpha$. 
\begin{enumerate}
    \item If $S$ is module-finite over $R$ and $J\subseteq S$ is an ideal, prove that $S/J$ is module-finite over $R$.
    \item If $S$ is module-finite over $R$ and $T$ is module-finite over $S$, prove that $T$ is module-finite over $R$.
    \item If $S$ is module-finite over $R$, and $R\to T$ is a ring map, prove that $T\otimes_R S$ is module-finite over $T$.
    \item If $U\subset R$ is multiplicatively closed and $S$ is module-finite over $R$, prove that $U^{-1}S$ is module-finite over $U^{-1}R$.
\end{enumerate}

\item \textbf{Comparing the Two Notions in Examples.} Decide whether each algebra is algebra-finite and whether it is module-finite over the indicated base ring. Justify your answers.
\begin{multicols}{2}
\begin{enumerate}
    \item $\KK[x]$ over $\KK$.
    \item $\KK[x]/\langle x^2\rangle$ over $\KK$.
    \item $\ZZ[i]$ over $\ZZ$.
    \item $\ZZ[1/2]$ over $\ZZ$.
    \item $\KK[t]$ over $\KK[t^2]$.
    \item $\KK[t,t^{-1}]$ over $\KK[t]$.
\end{enumerate}
\end{multicols}

\end{enumerate}

\hrulefill

Let $S$ be an $R$-algebra and let $s\in S$. We say that $s$ is \emph{integral over $R$} if there exists a monic polynomial with coefficients in $R$ that $s$ satisfies. In other words, $s$ is integral over $R$ if there are elements $a_1,\ldots,a_n\in R$ such that
\[
s^n+a_1s^{n-1}+a_2s^{n-2}+\cdots+a_{n-1}s+a_n=0
\]
in $S$, where each $a_i$ acts through the structure map $R\to S$. The word \emph{monic} is essential. For example, every rational number satisfies some nonzero polynomial over $\ZZ$, but only the integers satisfy monic polynomial equations over $\ZZ$. The monic condition is exactly what allows us to solve for high powers of $s$ in terms of lower powers.

\hrulefill

\begin{enumerate}[leftmargin=*, resume=problems]

\item \textbf{First Examples of Integral Elements.} Let $S$ be an $R$-algebra.
\begin{enumerate}
    \item Show that every element in the image of $R\to S$ is integral over $R$.
    \item Show that every nilpotent element of $S$ is integral over $R$.
    \item Let $S=R[x]/\langle x^2-r\rangle$ for some $r\in R$. Show that the image of $x$ in $S$ is integral over $R$.
    \item Show that $i\in \ZZ[i]$ is integral over $\ZZ$.
    \item Show that $t\in \KK[t]$ is integral over $\KK[t^2]$.
    \item Show that $\frac{1}{2}(\sqrt{-3}+1)\in \QQ(\sqrt{-3})$ is integral over $\ZZ[\sqrt{-3}]$. 
    \item Show directly that $1/2\in \QQ$ is not integral over $\ZZ$. Hint: Suppose $(1/2)^n+a_1(1/2)^{n-1}+\cdots+a_n=0$ and clear denominators.
\end{enumerate}

\item \textbf{Integral Elements Give Finite Modules.} Let $S$ be an $R$-algebra and let $s\in S$.
\begin{enumerate}
    \item Suppose that the element $s\in S$ is integral over $R$, and 
    \[
    s^n+a_1s^{n-1}+\cdots+a_n=0
    \]
    with $a_i\in R$. Show that every power $s^m$ with $m\geq n$ is an $R$-linear combination of $1,s,\ldots,s^{n-1}$.
    \item Deduce that $R[s]$ is generated as an $R$-module by $1,s,\ldots,s^{n-1}$.
    \item Find a finite set of $R$-module generators for $S=\KK[t]$ over $R=\KK[t^2]$
    \item Find a finite set of $R$-module generators for $S=\ZZ[i]$ over $R=\ZZ$. 
\end{enumerate}

\item \textbf{A Fundamental Criterion.} Let $S$ be an $R$-algebra and let $s\in S$. Prove that the following are equivalent.
\begin{enumerate}
    \item $s$ is integral over $R$.
    \item $R[s]$ is module-finite over $R$.
    \item There exists a finitely generated $R$-submodule $M\subseteq S$ such that $1\in M$ and $sM\subseteq M$.
\end{enumerate}

\item \textbf{Integral Elements Form a Subring.} Let $R \to S$ be an $R$-algebra. Let $\overline{R}^{S}$ denote the set of elements of $S$ that are integral over $R$. When $R$ is an integral domain and $S$ is a field we call $\overline{R}^{S}$ the \emph{integral closure of $R$ in $S$}. The \emph{integral closure} or \emph{normalization} of a domain $R$ is the integral closure of $R$ in its field of fractions $S=\Frac(R)$. We denote the integral closure of $R$ by $\overline{R}$. 
\begin{enumerate}
    \item If $x\in S$ is integral over $R$, prove that $R[x]$ is module-finite over $R$.

    \item If $x,y\in S$ are integral over $R$, prove that $R[x,y]$ is module-finite over $R$. Hint: First show $R[x]$ is finite over $R$, then show $R[x,y]$ is finite over $R[x]$.

    \item Let $s\in R[x,y]$. Show that multiplication by $s$ defines an
    $R$-linear map
    \[
    \mu_s:R[x,y]\longrightarrow R[x,y],
    \qquad
    f \longmapsto sf.
    \]
    Use Cayley--Hamilton, or the determinant trick, to prove that $s$ is
    integral over $R$.

    \item Conclude that $x+y$, $xy$, and $-x$ are integral over $R$.

    \item Conclude that $\overline{R}^{S}$ is a subring of $S$ containingnthe image of $R$.


\end{enumerate}

\item \textbf{Computations with Integrality.} Let $\KK$ be a field.
\begin{enumerate}
    \item Prove that the elements of $\QQ$ integral over $\ZZ$ are exactly the integers. Hint: Write a rational number in lowest terms.
    \item Show that $\KK[t]$ is module-finite over $\KK[t^2,t^3]$. Find explicit module generators.
    \item Show that $\KK[x,y]/\langle y^2-x^3\rangle$ is module-finite over $\KK[x]$. Find explicit module generators.
    \item Let $S=\KK[x,x^{-1}]$, viewed as a $\KK[x]$-algebra. Show that $x^{-1}$ is not integral over $\KK[x]$.
    \item Deduce that $\KK[x,x^{-1}]$ is algebra-finite but not module-finite over $\KK[x]$.
\end{enumerate}

\item \textbf{UFDs Are Integrally Closed.} Let $R$ be a UFD with fraction field $S=\Frac(R)$.
\begin{enumerate}
    \item Explain how we can think of $S$ as an $R$-algebra? 
    \item Let $\alpha\in S$. Explain why we can write $\alpha=\frac{a}{b}$ with $a,b\in R$, $b\neq 0$, and $a$ and $b$ having no common irreducible factors. 
    \item Give an example of a domain $R$ that is \emph{not} a UFD for which part (b) is false. 

    \item Suppose $\alpha=a/b \in S$ is integral over $R$ and
    \[
    \left(\frac{a}{b}\right)^n
    +r_1\left(\frac{a}{b}\right)^{n-1}
    +\cdots
    +r_{n-1}\left(\frac{a}{b}\right)
    +r_n
    =0
    \]
     for $r_1,\ldots,r_n\in R$. By clearing denominators show that $b$ divides $a^{n}$. 

    \item Show that if $a$ and $b$ have no common irreducible factor and $b\mid a^n$, then $b$ is a unit. Hint: Use unique factorization.

    \item Conclude that $\alpha=a/b\in R$. Therefore every element of $K$ that is integral over $R$ already lies in $R$.

    \item Conclude that every UFD is integrally closed.
    
    \item Show that $R=\KK[t^2,t^3]$ is not integrally closed by finding an explicitly integral element in $S=\Frac(R)=\KK(t)$ not contained in $R$. 
    
    \item Show that $\ZZ[\sqrt{-3}]$ is not integrally closed. What does this mean about it being a UFD?
\end{enumerate}

\end{enumerate}

\hrulefill

An $R$-algebra $S$ is called an \emph{integral extension} of $R$ if every element of $S$ is integral over $R$. Said differently, the structure map $R\to S$ is integral if all elements of the target satisfy monic polynomial equations over the source. Integral extensions are the precise setting where algebra generation and module generation interact especially well. The exercises below prove the fundamental comparison theorem:
\[
\boxed{\text{module-finite}}\quad \Longleftrightarrow\quad \boxed{\text{algebra-finite and integral}}.
\]
More explicitly, the exercises that follow will guide us through the proof of the following theorem. 

\begin{theorem}
Let $R\to S$ be an $R$-algebra.  The  $R$-algebra $S$ is module-finite over $R$ if and only if it $S$ is algebra-finite as an $R$-algebra. 
\end{theorem}

Note you will often, especially in algebraic geometry, see this expressed as finite is equivalent to finite type and integral. However, again we try to avoid this language for now. 

\hrulefill

\begin{enumerate}[leftmargin=*, resume=problems]

\item \textbf{Module-Finite Extensions Are Integral.} Let $S$ be a module-finite $R$-algebra.
\begin{enumerate}
    \item Fix $s\in S$. Show that multiplication by $s$ defines an $R$-linear map
    \[
    \mu_s:S\longrightarrow S
    \qquad \text{given by} \qquad
    x\longmapsto sx.
    \]
    \item Apply the Cayley-Hamilton theorem (explain why you can apply it) to $\mu_{s}$ to deduce that $s$ is integral over $R$.
    \item Conclude that every module-finite $R$-algebra is integral over $R$.
\end{enumerate}

\item \textbf{Finite Type Integral Extensions Are Module-Finite.} Suppose $S=R[s_1,\ldots,s_n]$ and each $s_i$ is integral over $R$.
\begin{enumerate}
    \item Prove that $R[s_1]$ is module-finite over $R$.
    \item Prove that $R[s_1,s_2]$ is module-finite over $R[s_1]$, and hence module-finite over $R$.
    \item Continue by induction to prove that $S$ is module-finite over $R$.
    \item Deduce that if $S$ is algebra-finite and integral over $R$, then $S$ is module-finite over $R$.
\end{enumerate}

\item \textbf{The Comparison Theorem.} Let $S$ be an $R$-algebra. Prove that the following are equivalent.
\begin{enumerate}
    \item $S$ is module-finite over $R$.
    \item $S$ is algebra-finite over $R$ and integral over $R$.
\end{enumerate}
Then revisit Problem 9 and explain how the theorem classifies each example.

\item \textbf{Permanence of Integral Extensions.} Let $R\to S\to T$ be maps of rings.
\begin{enumerate}
    \item Prove that every quotient $R/I$ is integral over $R$.
    \item Suppose $S$ is integral over $R$ and $T$ is integral over $S$. Prove that $T$ is integral over $R$. Hint: If $t\in T$ satisfies a monic equation over $S$, first adjoin to $R$ the finitely many coefficients appearing in that equation.
    \item If $S$ is integral over $R$ and $J\subseteq S$ is an ideal, prove that $S/J$ is integral over $R/(J\cap R)$.
    \item If $S$ is integral over $R$ and $U\subset R$ is multiplicatively closed, prove that $U^{-1}S$ is integral over $U^{-1}R$.
    \item More generally, if $S$ is integral over $R$ and $R\to R'$ is any ring map, prove that $R'\otimes_R S$ is integral over $R'$.
\end{enumerate}

\item \textbf{A First Consequence for Maximal Ideals.} Let $A\subseteq B$ be an integral extension of domains.
\begin{enumerate}
    \item Suppose $B$ is a field. Prove that $A$ is a field. Hint: If $0\neq a\in A$, then $a^{-1}\in B$ is integral over $A$; write a monic equation for $a^{-1}$ and multiply by a suitable power of $a$.
    \item Let $R\to S$ be an integral ring map and let $\fq\in\Spec(S)$ with contraction $\fp=\fq\cap R$. Show that $S/\fq$ is integral over $R/\fp$.
    \item Deduce that if $\fq$ is maximal, then $\fp$ is maximal.
    \item Prove the converse: if $\fp$ is maximal, then $\fq$ is maximal. Hint: $S/\fq$ is an integral domain integral over the field $R/\fp$.
    \item Conclude that under an integral ring map, a prime ideal upstairs is maximal if and only if its contraction downstairs is maximal.
\end{enumerate}

\item \textbf{More Examples of Integral Extensions.} Decide whether each map is integral. If it is integral and algebra-finite, give module generators.
\begin{enumerate}
    \item $\ZZ\to \ZZ[i]$.
    \item $\KK[t^2]\to \KK[t]$.
    \item $\KK[t^2,t^3]\to \KK[t]$.
    \item $\KK[x]\to \KK[x,x^{-1}]$.
    \item $\KK[x]\to \KK[x,y]/\langle y^2-x^3\rangle$.
    \item $\KK[x^2,xy,y^2]\to \KK[x,y]$. Hint: Show that $x$ and $y$ are integral over the source.
\end{enumerate}

\end{enumerate}

\hrulefill

We now introduce a numerical invariant of a ring. A \emph{chain of prime ideals of length $d$} in a ring $R$ is a strictly increasing chain
\[
\fp_0\subsetneq \fp_1\subsetneq \cdots \subsetneq \fp_d
\]
of prime ideals. Thus the length counts the number of strict inclusions, not the number of prime ideals appearing in the chain. For a nonzero ring $R$, the \emph{Krull dimension} of $R$ is
\[
\dim(R)
\coloneqq
\sup\left\{d\in\ZZ_{\geq 0}\;\middle|\;\text{there exists a chain }\fp_0\subsetneq\cdots\subsetneq\fp_d\text{ in }\Spec(R)\right\}.
\]
If there are chains of arbitrarily large finite length, we write $\dim(R)=\infty$. The zero ring has no prime ideals; different authors use different conventions for its dimension, so we will avoid using the zero ring in basic dimension computations unless a convention is specified. We will eventually prove the following theorem, which will turn out to be key to computing the dimension of most rings.

\begin{theorem}
	If $\KK$ is a field then $\dim \KK[x_1,x_2,\ldots,x_n]=n$.
\end{theorem}

Proving that the dimension of the polynomial ring $\KK[x_1,x_2,\ldots,x_n]$ is at least $n$ is relatively easy, all one must do is find an explicit chain of prime ideals of length $n$. You will do this in the exercises below. Proving the equality in the theorem requires us to prove the other inequality, i.e. we must show there is no chain of prime ideals of length more than $n$. This turns out to be surprisingly difficult. For the remainder of the worksheet you may assume this theorem unless it trivializes an exercise. 

You will also prove that if $R$ is any ring then $\dim R[x] \geq \dim R + 1$, again by exhibiting an explicit chain of prime ideals. Intuition tells us that this should be an equality, however, this turns out to be false unless we put additional hypotheses on $R$, namely that $R$ be Noetherian. This highlights a somewhat tedious fact about dimension theory in commutative algebra: unless one assumes some form of finiteness condition (Noetherian, finitely generated $\KK$-algebra, integral, etc.) there are many pathologies. 

Using the correspondence between prime ideals and irreducible closed subsets of $\Spec(R$), thee intuition is that longer chains of prime ideals correspond to longer chains of irreducible closed subsets in $\Spec(R)$, ordered in the opposite direction. 

\hrulefill

\begin{enumerate}[leftmargin=*, resume=problems]

\item \textbf{Unpacking the Definition.} Let $R$ be a nonzero ring.
\begin{enumerate}
    \item Show that $\dim(R)=0$ if and only if every prime ideal of $R$ is maximal.
    \item Show that every field has Krull dimension $0$.
    \item Show that if $R$ is an integral domain that is not a field, then $\dim(R)\geq 1$.
    \item Let $I\subseteq R$ be an ideal. Use the correspondence between prime ideals of $R/I$ and prime ideals of $R$ containing $I$ to describe $\dim(R/I)$ in terms of chains of prime ideals in $R$.
    \item Prove that $\dim(R)=\dim(R_{\red})$, where $R_{\red}=R/\Nil(R)$.
\end{enumerate}

\item \textbf{Dimension Zero Examples.} Compute the Krull dimension of each ring.
\begin{multicols}{2}
\begin{enumerate}
    \item $\KK$.
    \item $\KK\times \KK$.
    \item $\KK[\epsilon]/\langle \epsilon^2\rangle$.
    \item $\ZZ/12\ZZ$.
    \item $\KK[x]/\langle x^n\rangle$ for $n\geq 1$.
    \item $\KK[x,y]/\langle x,y\rangle$.
\end{enumerate}
\end{multicols}

\item \textbf{Dimension One Examples.} 
\begin{enumerate}
    \item Prove that $\dim(\ZZ)=1$.
    \item More generally, prove that every principal ideal domain that is not a field has Krull dimension $1$.
    \item Prove that $\dim(\KK[x])=1$.
    \item Compute $\dim(\KK[x]/\langle f\rangle)$ for a nonzero nonunit polynomial $f\in \KK[x]$.
    \item Compute $\dim(\ZZ[1/n])$ for an integer $n\geq 2$.
\end{enumerate}

\item \textbf{Building Longer Chains.} Let $\KK$ be a field.
\begin{enumerate}
    \item Prove that
    \[
    \langle 0\rangle
    \subsetneq
    \langle x_1\rangle
    \subsetneq
    \langle x_1,x_2\rangle
    \subsetneq
    \cdots
    \subsetneq
    \langle x_1,\ldots,x_n\rangle
    \]
    is a chain of prime ideals in $\KK[x_1,\ldots,x_n]$.
    \item Deduce that $\dim(\KK[x_1,\ldots,x_n])\geq n$.
    \item Let $R$ be any nonzero ring. Show that $\dim(R[x])\geq \dim(R)+1$ whenever $\dim(R)$ is finite. Hint: Extend a chain in $R$ to $R[x]$, then add the ideal generated by the final prime together with $x$.
\end{enumerate}

\item \textbf{Dimension and Quotients.} Let $\KK$ be a field.
\begin{enumerate}
    \item Compute $\dim(\KK[x,y]/\langle xy\rangle)$. Hint: Every prime ideal containing $xy$ contains $x$ or $y$.
    \item Compute $\dim(\KK[x,y]/\langle x^2,xy\rangle)$. Hint: First pass to the reduced ring.
    \item Compute $\dim(\KK[x,y]/\langle y-x^2\rangle)$.
    \item Find a chain of prime ideals of length $2$ in $\KK[x,y,z]/\langle xy\rangle$.
    \item Assuming $\dim(\KK[x_1,\ldots,x_n])=n$, compute the dimension of $\KK[x,y,z]/\langle z-y^2+x^3\rangle$.
\end{enumerate}

\item \textbf{Local Dimension.} Let $\fp\in\Spec(R)$.
\begin{enumerate}
    \item Recall that prime ideals of $R_{\fp}$ correspond to prime ideals $\fq\subseteq R$ with $\fq\subseteq \fp$. Use this to describe $\dim(R_{\fp})$ in terms of chains of prime ideals in $R$ ending inside $\fp$.
    \item Compute $\dim(\KK[x]_{\langle x\rangle})$.
    \item Compute $\dim(\ZZ_{\langle p\rangle})$ for a prime number $p$.
    \item Find a chain showing that $\dim(\KK[x,y]_{\langle x,y\rangle})\geq 2$.
    \item Assuming $\dim(\KK[x,y])=2$, prove that $\dim(\KK[x,y]_{\langle x,y\rangle})=2$.
\end{enumerate}

\end{enumerate}

\end{document}
